\subsubsection{延拓定理}\label{subsec:sobolev-extension-theorems}

现在讨论能否将 $H^k(\Omega)$ 中的函数延拓为 $H^k(\mathbb{R}^m)$ 中的函数. 先指出一个显然的延拓结果.

\begin{Lemma}
\label{lem:zero-boundary-sobolev-zero-extension}
设 $u\in H_0^k(\Omega)$. 那么通过在 $\overline\Omega$ 外定义为零, 可将 $u$ 延拓为 $H^k(\mathbb{R}^m)$ 中的函数, 即对每个 $\phi\in\mathcal{D}(\mathbb{R}^m)$, 令 $(u,\phi)=\int_{\Omega}\overline{u(\bm{x})}\phi(\bm{x})\,\mathrm{d}\bm{x}$.
\end{Lemma}

\begingroup
\postdisplaypenalty=10000
\renewcommand{\endProof}{\par\nobreak\hfill$\blacksquare$\par\vspace{0.1cm}}
\begin{Proof}
取在 $H^k(\Omega)$ 中收敛到 $u$ 的序列 $u_n\in\mathcal{D}(\Omega)$. 那么 $u_n$ 也在 $H^k(\mathbb{R}^m)$ 中收敛.
\end{Proof}
\endgroup

接着考察半空间的情形. 令 $\mathbb{R}_+^m=\{x_m>0\}$. 有以下结果.

\begin{Theorem}
\label{thm:halfspace-sobolev-bounded-extension}
设 $k\geq0$ 为整数. 那么存在有界线性映射 $E:H^k(\mathbb{R}_+^m)\to H^k(\mathbb{R}^m)$, 使对每个 $u\in H^k(\mathbb{R}_+^m)$ 有 $(Eu)|_{\mathbb{R}_+^m}=u$. 此外, 对任意给定的 $K\in\mathbb{N}$, 可对 $0\leq k\leq K$ 选取与 $k$ 无关的 $E$.
\end{Theorem}

\begingroup
\postdisplaypenalty=10000
\renewcommand{\endProof}{\par\nobreak\hfill$\blacksquare$\par\vspace{0.1cm}}
\begin{Proof}
先考察如何延拓光滑函数. 连续函数可通过偶延拓成为整个空间上的连续函数, 即对 $x_m<0$, 令 $u(\bm{x}^{\prime},x_m)=u(\bm{x}^{\prime},-x_m)$. 这里及下文中, $\bm{x}^{\prime}$ 表示 $(x_1,x_2,\ldots,x_{m-1})$. 然而, 偶延拓一般会导致导数不连续. 若对 $x_m<0$ 令 $u(\bm{x}^{\prime},x_m)=4u(\bm{x}^{\prime},-x_m/2)-3u(\bm{x}^{\prime},-x_m)$, 则当 $u$ 光滑时, $u$ 和 $\partial u/\partial x_m$ 都在跨过 $x_m=0$ 时连续. 一般地, 对 $x_m<0$ 作如下设定:
\begin{equation}\label{eq:halfspace-extension-reflected-linear-combination}
u(\bm{x}^{\prime},x_m)=\sum_{j=0}^K\alpha_j u\left(\bm{x}^{\prime},-\dfrac{x_m}{j+1}\right)
\end{equation}
希望选择 $\alpha_j$, 使对光滑的 $u$, 直到 $K$ 阶的导数在 $x_m=0$ 两侧匹配. 这得到方程组
\begin{equation}\label{eq:halfspace-extension-vandermonde-matching-system}
\sum_{j=0}^K\left(-\dfrac{1}{j+1}\right)^i\alpha_j=1,\quad i=0,1,\ldots,K.
\end{equation}
该方程组的矩阵称为 Vandermonde 矩阵, 已知它非奇异 (Problem~\ref{ex:vandermonde-matrix-nonsingularity}). 因而 \eqref{eq:halfspace-extension-vandermonde-matching-system} 唯一可解.

容易看出, 若 $u\in C^k(\overline{\mathbb{R}_+^m})$ $(k\leq K)$ 且具有有界支集, 则 $Eu\in C^k(\mathbb{R}^m)$ 且具有有界支集, 并对某个与 $u$ 无关的常数 $C$ 有 $\|Eu\|_{k,2}\leq C\|u\|_{k,2}$. 若能证明 $C^k(\overline{\mathbb{R}_+^m})$ 中具有有界支集的函数在 $H^k(\mathbb{R}_+^m)$ 中稠密, 就完成了证明.
\end{Proof}
\endgroup

\begin{Lemma}
\label{lem:halfspace-restricted-test-function-density}
设 $u\in H^k(\mathbb{R}_+^m)$. 那么存在序列 $\phi_n\in\mathcal{D}(\mathbb{R}^m)$, 使 $\phi_n$ 在 $\mathbb{R}_+^m$ 上的限制收敛到 $u$.
\end{Lemma}

\begingroup
\postdisplaypenalty=10000
\renewcommand{\endProof}{\par\nobreak\hfill$\blacksquare$\par\vspace{0.1cm}}
\begin{Proof}
由于已知 $\mathcal{D}(\mathbb{R}^m)$ 在 $H^k(\mathbb{R}^m)$ 中稠密, 只需证明每个 $u\in H^k(\mathbb{R}_+^m)$ 都可由 $H^k(\mathbb{R}^m)$ 中函数的限制逼近. 此外, 已知有界支集函数稠密; 因而假设 $u$ 具有有界支集. 对 $\bm{x}\in\mathbb{R}_+^m$, 令 $u_n(\bm{x}^{\prime},x_m)=u(\bm{x}^{\prime},x_m+1/n)$. 那么 $\|u_n\|_2\leq\|u\|_2$, 且对 $u\in\mathcal{D}(\mathbb{R}_+^m)$ 有 $u_n\to u$. 由稠密性, 对每个 $u\in L^2(\mathbb{R}_+^m)$, 有 $u_n\to u$ 在 $L^2$ 中成立. 此外, 由于 $D^\alpha u_n(\bm{x}^{\prime},x_m)=D^\alpha u(\bm{x}^{\prime},x_m+1/n)$, 同样的论证表明 $u_n\to u$ 在 $H^k(\mathbb{R}_+^m)$ 中成立. 断言 $u_n$ 可以延拓. 首先注意 $u_n$ 实际上在 $x_m>-1/n$ 上有定义. 令
\begin{equation}\label{eq:halfspace-density-shifted-zero-extension-distribution}
(v_n,\phi)=\int_{x_m>-1/n}\overline{u_n(\bm{x})}\phi(\bm{x})\,\mathrm{d}\bm{x}
\end{equation}
为 $u_n$ 的零延拓. 那么显然, 当 $x_m<-1/n$ 时 $D^\alpha v_n=0$, 当 $x_m>-1/n$ 时 $D^\alpha v_n=D^\alpha u_n$. 然而, 不能断言 $D^\alpha v_n$ 是 $D^\alpha u_n$ 的零延拓. 可以断言的是
\begin{equation}\label{eq:halfspace-density-zero-extension-boundary-distribution}
(D^\alpha v_n,\phi)=\int_{x_m>-1/n}\overline{D^\alpha u_n(\bm{x})}\phi(\bm{x})\,\mathrm{d}\bm{x}+(g_{n\alpha},\phi),
\end{equation}
其中 $g_{n\alpha}$ 为支集有界且包含于平面 $x_m=-1/n$ 的分布. 由 Lemma~\ref{lem:distribution-local-finite-order}, $g_{n\alpha}$ 具有有限阶. 现在设 $\psi_n(x_m)$ 为 $C^{\infty}$ 函数, 在 $x_m\leq-1/n$ 时为零, 在 $x_m\geq0$ 时等于 $1$. 那么 $g_{n\alpha}\psi_n$ 为零, 因而 $w_n=v_n\psi_n$ 属于 $H^k(\mathbb{R}^m)$. 显然, $w_n$ 在 $\mathbb{R}_+^m$ 上的限制为 $u_n$. 证明完毕.
\end{Proof}
\endgroup
